% ID: SEL-D0080-OB-E03
% ESTADO: cotejada. Segunda pasada visual y borrador compilado revisados.
% Fuente DOC: Selectividad/sel_2012_matematicas_aplicadas/2 MATEMÁT. APLIC. CC. SS/EXAMEN 5-ANDALUCÍA.doc
% SHA-256 original: 87b729cdd4632386690a9390bdfa17dd7d5c902a2cf7e5801feea45e2b181969
% Página convertida: 2; opción B; ejercicio 3.

Se sabe que el $90\%$ de los estudiantes del último curso de una Universidad está preocupado por sus posibilidades de encontrar trabajo, el $30\%$ está preocupado por sus notas y el $25\%$ por ambas cosas.
\begin{enumerate}[label=\alph*)]
\item \textbf{(1.5 puntos)} Si hay 400 alumnos matriculados en el último curso de dicha Universidad, \textquestiondown{}cuántos de ellos no están preocupados por ninguna de las dos cosas?
\item \textbf{(1 punto)} Si un alumno del último curso, elegido al azar, no está preocupado por encontrar trabajo, \textquestiondown{}cuál es la probabilidad de que esté preocupado por sus notas?
\end{enumerate}
